% =====================================================================
% Mục III: pipeline chung và ba cấp độ (03-1, 03-2, 03-3). Ablation của từng cấp nằm cuối mục con của cấp đó;
% kết quả chính trên test ở Mục IV.
\section{Pipeline}
\label{sec:pipeline}

Cả ba cấp độ dùng chung một pipeline, đúng quy trình gợi ý trong đề (Hình~\ref{fig:pipeline}).
\begin{enumerate}
\item CLIP ViT-B/32 (trọng số OpenAI, đóng băng) mã hoá câu truy vấn và lấy 50 ảnh có cosine lớn nhất (FAISS~\cite{johnson2019faiss}
ở VG và demo; phép nhân ma trận trực tiếp ở MKG-W vì pool chỉ có 892 ảnh).
\item Với từng ứng viên, một mô-đun đồ thị tính điểm đồ thị. Ở VG, đó là điểm khớp giữa đồ thị câu và scene graph của ảnh.
Ở MKG-W, đó là quan hệ giữa thực thể ứng viên và các thực thể được nhắc trong câu.
\item Điểm cuối kết hợp hai điểm theo dạng $\alpha\cdot\text{CLIP} + (1-\alpha)\cdot\text{đồ thị}$ của đề. Mô hình cấp~1
và MKG-W chuẩn hoá từng điểm bằng z-score trong 50 ứng viên của truy vấn; mô hình chính của VG cộng điểm thô
(Mục~\ref{sec:method-main}). Ứng viên được sắp lại theo điểm cuối.
\item Lời giải thích của mỗi ứng viên là chính các đại lượng đã cộng vào điểm đồ thị. Không có bước tìm đường đi nào chạy
sau khi đã có thứ hạng (Mục~\ref{sec:method-expl} và~\ref{sec:l2-expl}).
\end{enumerate}
Vì chỉ xếp lại 50 ứng viên, giới hạn trên của mọi mô hình là recall@50 của CLIP. Trên test VG, chỉ số này là 96,66\%.
Mục~\ref{sec:method-l1}--\ref{sec:method-l3} mô tả từng cấp độ; mỗi mục kết thúc bằng các ablation dùng để kiểm tra
thành phần của cấp đó. Dữ liệu, cách chia tập và kết quả chính trên test ở Mục~\ref{sec:exp}.

\begin{figure}[t]
\centering
\begin{tikzpicture}[font=\scriptsize, node distance=3mm and 4mm,
  box/.style={draw, rounded corners=2pt, align=center, minimum height=7mm, inner sep=2pt},
  arr/.style={-{Stealth[length=1.6mm]}, thick, cGray}]
\node[box, fill=cClip!12] (q) {Truy vấn\\(câu)};
\node[box, fill=cClip!12, right=of q] (clip) {CLIP\\đóng băng};
\node[box, fill=cClip!12, right=of clip] (top) {Top-50\\ứng viên};
\node[box, fill=cObj!15, below=6mm of q] (qg) {Đồ thị câu /\\thực thể nhắc tới};
\node[box, fill=cObj!15, right=of qg] (gs) {Điểm đồ thị\\(GATv2 / KG)};
\node[box, fill=cTri!12, right=of gs] (ig) {Scene graph ảnh /\\đồ thị tri thức};
\node[box, fill=yellow!20, below=6mm of gs] (fuse) {$\alpha\cdot\text{CLIP} + (1-\alpha)\cdot\text{đồ thị}$};
\node[box, right=of fuse] (out) {Thứ hạng\\+ giải thích};
\draw[arr] (q) -- (clip); \draw[arr] (clip) -- (top);
\draw[arr] (q) -- (qg); \draw[arr] (qg) -- (gs); \draw[arr] (ig) -- (gs);
\draw[arr] (top) -- (ig);
\draw[arr] (gs) -- (fuse); \draw[arr] (top.south) |- ($(fuse.north east)+(0,0.5mm)$);
\draw[arr] (fuse) -- (out);
\end{tikzpicture}
\caption{Pipeline chung. Phần màu xanh là CLIP đóng băng. Phần đồ thị chỉ chấm lại 50 ứng viên của CLIP. Lời giải thích
là các đại lượng của chính điểm đồ thị.}
\label{fig:pipeline}
\end{figure}
